\section*{第 \ref{ch:b1:lipids} 章——脂质}

\begin{solution}{exo:b1:lipids:1}
脂质是一类不溶于水、而溶于非极性溶剂的
生物分子。\omterm{def:b1:lipids:triglyceride}{脂肪}、\omterm{def:b1:lipids:amphiphilic}{磷脂}、\omterm{def:b1:lipids:amphiphilic}{固醇}和类胡萝卜素
并不共享同一骨架，它们共同的只是这种与水的相互作用方式，
正是这一点使它们承担共同的作用（储备、成膜、构成膜）。
\end{solution}

\begin{solution}{exo:b1:lipids:2}
C18:2 $\Delta^{9,12}$。最后一个双键始于距羧基第 12 位碳，
即距甲基端第 6 位碳，因此为 ω-6。
\end{solution}

\begin{solution}{exo:b1:lipids:3}
\omterm{def:b1:lipids:triglyceride}{甘油三酯}没有极性头部：它完全疏水，
在水中以本体相的形式被排斥，聚集成液滴。\omterm{def:b1:lipids:amphiphilic}{磷脂}则是\omterm{def:b1:lipids:amphiphilic}{两亲的}：
其极性头部必须留在水中，而其尾部必须离开水，
只有两个分子厚的片层才能同时满足这一要求。
\end{solution}

\begin{solution}{exo:b1:lipids:4}
约 \qty{41}{\celsius}、\qty{-5}{\celsius} 和 \qty{20}{\celsius}
（\omterm{def:b1:lipids:amphiphilic}{胆固醇}曲线没有尖锐的转变；其升高的一半位于
\qty{20}{\celsius} 附近）。在 \qty{10}{\celsius} 下的细菌需要不饱和组成，
使膜在远低于其生长温度时仍保持流动。
\end{solution}

\begin{solution}{exo:b1:lipids:5}
$15\,000\times 38 = \qty{570}{MJ}$；$570/8 = 71$ 天。糖原：干重
$570\,000/17 = \qty{33.5}{kg}$，含水后为 \qty{134}{kg}。
\end{solution}

\begin{solution}{exo:b1:lipids:6}
C22:0 $>$ C18:0 $>$ C14:0 $>$ C18:1 $>$ C18:3。在\omterm{def:b1:lipids:fattyacid}{饱和脂肪酸}中，
链越长，范德瓦耳斯接触越多；一个\emph{顺式}双键所破坏的堆积
多于额外四个碳所增加的堆积（C18:1 的熔点
低于 C14:0）；每增加一个双键，熔点又进一步降低。
\end{solution}

\begin{solution}{exo:b1:lipids:7}
$2\times\qty{140}{\micro m^2}/\qty{0.6}{nm^2} = 2\times 1.4\times
10^{-10}/6\times 10^{-19} = \num{4.7e8}$ 个分子；质量
$\num{4.7e8}\times 750/\num{6e23} = \qty{5.8e-13}{g} = \qty{0.58}{pg}$。
\end{solution}

\begin{solution}{exo:b1:lipids:8}
体积 \qty{1.11}{mL} $= \qty{1.11e-6}{m^3}$；球形表面积
$= 6V/d = 6\times 1.11\times 10^{-6}/10^{-6} = \qty{6.7}{m^2}$。若将该体积
汇聚成一滴，则 $d = \qty{1.28}{cm}$，表面积 \qty{5.2}{cm^2}：
乳液的界面面积增加了约一万三千倍。\omterm{def:b1:lipids:triglyceride}{脂肪}酶
溶于水，只在界面上起作用；其反应速率
与界面面积成正比。
\end{solution}

\begin{solution}{exo:b1:lipids:9}
高于 $T_m$ 时，其刚性环阻碍邻近链的运动，
降低\omterm{def:b1:lipids:fluidity}{流动性}；低于 $T_m$ 时，其体积庞大，阻止链
堆积成有序的凝胶，并阻止其凝固。由于无法发生
协同堆积，因此不存在使整个双分子层同时改变状态的温度：
其转变是分散进行的。
\end{solution}

\begin{solution}{exo:b1:lipids:10}
全饱和的链：在 \qty{25}{\celsius} 时，膜已
接近其转变点且较为僵硬；在 \qty{5}{\celsius} 时则发生凝胶化。
凝胶化的膜使蛋白质停止工作，并在复温或
受到机械应力时破裂、渗漏：\omterm{def:b1:cell-unit-of-life:cell}{细胞}丧失内容物，
\omterm{def:b1:body-plans-tissues:tissue}{组织}随之死亡（冷害）。野生型在秋季使其脂质去饱和，
并在 \qty{5}{\celsius} 时仍保持流动。
\end{solution}

\begin{solution}{exo:b1:lipids:11}
每千克\omterm{def:b1:lipids:triglyceride}{脂肪}产生 \qty{1.07}{kg} 水。氧气：\qty{2000}{L}；
若以 \qty{5}{\%} 的提取率从含氧 \qty{21}{\%} 的空气中获取，所需空气
$2000/(0.21\times 0.05) = \qty{190}{m^3}$；以
\qty{30}{g/m^3}计，失水 \qty{5.7}{kg}。骆驼呼吸
失去的水比\omterm{def:b1:lipids:triglyceride}{脂肪}产生的水多五倍：驼峰是能量储存，
而非水的储存。
\end{solution}

\begin{solution}{exo:b1:lipids:12}
对于双分子层，\omterm{prop:b1:water-small-molecules:properties}{疏水效应}使其重新封闭、闭合成
囊泡，并在无须能量输入的情况下接纳新脂质：组装是
自发的。但\omterm{def:b1:cell-unit-of-life:cell}{细胞}借助镶嵌于已有膜（\omterm{def:b1:eukaryotic-cell:endomembrane}{内质网}）中的酶来合成脂质，
这些酶就地将脂质插入其中；新膜通过旧膜的扩展而生长，
并在\omterm{def:b1:cell-unit-of-life:cell}{细胞}分裂时被分配，因此
每个\omterm{def:b1:cell-unit-of-life:cell}{细胞}的每一层膜都源自
上一个\omterm{def:b1:cell-unit-of-life:cell}{细胞}的膜。自组装解释的是片层的物理学，
而非其在\omterm{def:b1:cell-unit-of-life:cell}{细胞}中的起源。
\end{solution}

\begin{solution}{pb:b1:lipids:1}
\textbf{1.} C16:0（0），C18:1 $\Delta^9$（1），C22:6 $\Delta^{4,7,10,13,16,19}$
（6）。
\textbf{2.} 温暖时：$0.40\times 0 + 0.45\times 1 + 0.15\times 6 = 1.35$；
寒冷时：$0 + 0.40 + 2.10 = 2.50$ 个双键（每条链）。
\textbf{3.} 每个\emph{顺式}键使链弯折，并阻止其与邻链紧密
堆积；范德瓦耳斯接触减少，
使链失序所需的热量减少，$T_m$ 降低。
\textbf{4.} 每条链多 $1.15$ 个双键，其\omterm{def:b1:lipids:fluidity}{相变温度}
约低 \qty{45}{\celsius}。
\textbf{5.} 在凝胶中，脂质无法移动：\omterm{def:b1:membranes-transport:transporters}{载体蛋白}无法改变
构象，泵停止工作，\omterm{def:b1:membranes-transport:transporters}{通道蛋白}卡住；离子梯度消散，神经
传导失效，鱼体瘫痪，并死于
经过冷适应的鱼本可耐受的寒冷。
\textbf{6.} 去饱和酶，它们将双键引入脂酰
链中；它们是\omterm{def:b1:eukaryotic-cell:endomembrane}{内质网}的\omterm{def:b1:membranes-transport:proteins}{整合蛋白}，脂质
正是在此合成。
\textbf{7.} 在两种情况下，它都起到缓冲\omterm{def:b1:lipids:fluidity}{流动性}的作用：使
非常不饱和的冷适应膜在其较低的 $T_m$ 之上变得僵硬，并使温暖的
膜在凉爽的日子里不致凝胶化。
\textbf{8.} $30\,000\times 38 = \qty{1140}{MJ}$。
\textbf{9.} $1140/60 = 19$ 天。
\textbf{10.} 干重：$\num{1.14e6}/17 = \qty{67}{kg}$；含水后：
\qty{268}{kg}。
\textbf{11.} $268 - 30 = \qty{238}{kg}$，几乎是骆驼体重的一半。
\textbf{12.} 若在全身皮肤下都有一层\omterm{def:b1:lipids:triglyceride}{脂肪}，身体便会被隔热，
难以在沙漠中散热；而把\omterm{def:b1:lipids:triglyceride}{脂肪}集中在一个驼峰中，
则身体其余部分的皮肤仍可散热。
\textbf{13.} $1.07\times 30 = \qty{32}{kg}$ 水。
\textbf{14.} $2.0\times 30\,000 = \qty{60000}{L} = \qty{60}{m^3}$ 的
氧气。
\textbf{15.} $60/(0.21\times 0.05) = \qty{5700}{m^3}$ 空气。
\textbf{16.} $5700\times 30 = \qty{171}{kg}$ 水。
\textbf{17.} 失去的水是获得水量的五倍：氧化驼峰要消耗
水分。驼峰是一种能量储存，使骆驼可以不进食
也能生存；其节水能力来自肾脏、对脱水的耐受力，
以及允许体温升高的能力。
\textbf{18.} $500\times 3.5\times 6 = \qty{10.5}{MJ}$，以热量形式储存并于夜间释放；若以蒸发方式将其散去，则每天将损失 $10\,500/2.4 = \qty{4.4}{kg}$ 水。
\textbf{19.} 膜只能拉伸几个百分点，因此要使红\omterm{def:b1:cell-unit-of-life:cell}{细胞}
体积加倍，就需要大量过剩的膜，并将其折叠成静止时的
双凹形态；还需要一个\omterm{def:b1:eukaryotic-cell:cytoskeleton}{细胞骨架}，使膜能够展开而不撕裂：
骆驼的红\omterm{def:b1:cell-unit-of-life:cell}{细胞}呈椭圆形，单位体积所含的膜
比人类的更多。
\textbf{20.} $2\times 80\times 10^{-12}/0.65\times 10^{-18} = \num{2.5e8}$。
\textbf{21.} $\num{2.5e8}\times\num{8e12}\times 40 = \num{7.9e22}$。
\textbf{22.} $\num{7.9e22}\times 750/\num{6e23} = \qty{99}{g}$。
\textbf{23.} 侧向扩散使头部始终在水中、尾部始终在油中，
每一步都无需付出代价；而翻转要求极性头部
穿过疏水核心，这是一道很高的能垒，只有热
运动才能每周一次地跨越它。
\textbf{24.} 两个小叶之间的差异只能以与
脂质翻转相同的速率衰减；每个分子每周约翻转一次。
而由能够将特定脂质转运过膜的酶（翻转酶）所建立的
不对称性，则可无限期地抵抗自发翻转而保持。
\textbf{25.} 约节省 \qty{240}{kg}——几乎是其体重的一半
——其做法是携带 \qty{30}{kg} \omterm{def:b1:lipids:triglyceride}{脂肪}，而非 \qty{270}{kg} 含水糖原；
而驼峰并非水的储存：氧化它时经呼吸损失的水，
是其所产生水量的五倍。
\end{solution}
